\scanpage{13}
\deleted{2. Soit $X$ un espace, $Y$ un sous-espace fermé.}
\deleted{$(\unsure{6})$ Soit $\bar T$}

\uline{NB} Les $\bar T_\lambda$ $(\lambda\in[0,1])$ forment
une famille croissante de parties fermées de
$\bar T$, les $\mathring T_\lambda\overset{\mathrm{déf}}=
\lambda\mathring T\deleted{\unsure{_\lambda}}$ $(\lambda\in]0,1])$
une famille croissante de parties ouvertes. Si $\lambda\in]0,1[$, alors
$\mathring T_\lambda$ est bel et bien l'intérieur \inserted{topologique}
de $\bar T_\lambda$ dans $\bar T$,
mais bien sûr $\mathring T_1=\mathring T$ n'est pas l'intérieur
\inserted{top.} de $\bar T_1$ dans
l'espace topologique $\bar T_1=\bar T$, i.e. on n'a pas
$\mathring T\deleted{\unsure{_1}}=\bar T$ sauf
si $\dot T=\varnothing$ i.e. $Y=\bar T$.

Posons, pour $\lambda\in]0,1[$
\[
(\unsure{6})\qquad
\left\{\begin{aligned}
\bar T\setminus\mathring T_\lambda
&\overset{\mathrm{déf}}=\deleted{\unsure{E}}
\mathrm{Ep}^{(T)}_{1-\lambda=\varepsilon}
\quad\text{``épaississement d'indice $\varepsilon$''}\\
&\hspace{4em}\text{du bord $\dot T$ de $\bar T$}\\
\dot{\mathrm{Ep}}_\varepsilon^{(T)}
&\overset{\mathrm{déf}}=\dot T\cup(1-\varepsilon).\dot T
\end{aligned}\right.
\]
\uline{NB} La frontière de $E_\varepsilon$ dans $\bar T$ est
$(1-\varepsilon)\dot T$, et $E_\varepsilon$
est un voisinage de $\dot T$. On a
\[
(7)\qquad\mathrm{Ep}_\varepsilon^{(T)}
\simeq[\underbrace{1-\varepsilon}_{\lambda},1]\times\dot T
\]
\margin{\uline{NB} Si $Y,\dot T$ \unsure{compacts}, \deleted{\unsure{les}}
alors $\bar T$ compact, et les $\bar T_\varepsilon$
$(\varepsilon\in]0,1\unsure{[|]})$ forment
un syst. fond. de voisinages de $Y$ dans $\bar T$,
\deleted{fermés}\inserted{\unsure{ouverts}}.
(\unsure{ainsi que voisinages} $\unsure{\bar T_\varepsilon}$ tubulaires absolus de $Y$)}

\textcircled{2} Soit $Y\hookrightarrow X$ une inclusion fermée d'espaces
topologiques. Un \uline{voisinage conique} de \uline{$Y$ dans} $X$
est un voisinage $\bar T$ de $Y$ dans $X$
muni d'une opération (1) $I\times\bar T\longrightarrow\bar T$
\inserted{faisant de $\bar T$ un vois. \deleted{tub.}\inserted{coniq.}
absolu de $Y$.} \deleted{satisfaisant les conditions id.}
Alors les
$\lambda\bar T=\bar T_\lambda$ $(\lambda\in]0,1])$
forment d'autres voisinages
tubulaires de $Y$ dans $X$ (famille croissante) et
on voit que pour $\lambda\in]0,1[$, on a
\[
(8)\qquad\dot T_\lambda=\lambda\dot T
\quad\text{est la frontière de }\bar T_\lambda=\lambda\bar T\text{ dans }X.
\]
\margin{12}
